\thispagestyle{plain}
\section*{Acknowledgements}
\addcontentsline{toc}{chapter}{Acknowledgements}
\markboth{Acknowledgements}{}

\textit{To my family.}

\vspace{0.5cm}

Adóitase dicir que unha persoa ten dúas familias: aquela na que nace e a que elixe. Botando a vista atrás aos anos dedicados a esta tese doutoral, síntome profundamente afortunado de contar co apoio de ambas.

No que respecta á miña familia elixida no mundo académico, gustaríame expresar primeiro o meu agradecemento aos membros do grupo experimental Granja-Montenegro: Román, MVP, Zeko, Alba, Bayón, Lucía, Patri, Vicky, Amaia, Kaddy, Martín, Selene... (e aos profesores Granja e Amorín!), especialmente a Serantes, compañeiro dende o inicio da carreira, e a Alicia polas súas reivindicacións sobre o mundo académico (poderiamos formar un sindicato!). Con todo, a miña maior alegría foi, sen dúbida, o grupo computacional. Estou inmensamente agradecido polos amplos coñecementos de Fabián, a franqueza de Fonso e a súa capacidade para ir directo ao grao, e as singulares interpretacións filosóficas de Xandre. Tamén dou as grazas a Paula pola calma que transmite e a súa filosofía de gozar da vida; a Marta pola súa motivación constante; a Anxo pola súa tranquilidade e as súas infinitas anécdotas, e a David pola súa contaxiosa alegría e experiencia. Unha mención moi especial merecen Martín Calvelo e Pablo pola súa orientación e ensinanzas durante os meus primeiros pasos nesta viaxe. E, por suposto, a un sopro de aire fresco: Dani e Uxía.

Como non, debo mencionar á miña man dereita ao longo desta tese, Alejandro Seco. Desde as nosas ideas e discusións ata os proxectos compartidos, e mesmo a convivencia durante unha estadía de investigación en Guarda, formamos un dúo mellor ca \textit{Mortadelo e Filemón}! Tamén teño unha gran débeda cos meus directores, Ángel e Rebeca. Grazas pola vosa confianza, pola vosa axuda na busca dun posdoutorado (aínda que non era a vosa responsabilidade) e polos vosos consellos máis alá do eido académico e os debates sobre o futuro da humanidade.

Finalmente, quero agradecer a acollida recibida nas dúas estadías realizadas, con Cecilia Clementi e o seu humor xa alemán... especialmente a Daniel Ramírez e a Aldo Pasos-Trejo, por seren os de fala española cos que mellor puiden colaborar. Tamén quero agradecer a Hugo A. L. Filipe e a Antonio, por esas tardes de viños POPA (non-lípido) en Guarda. E a Carme Rovira, por ensinarme que a Dinámica Molecular podería complicarse aínda moito máis cos PMF en QM/MM.

Unha mención especial para Yeray (a pesar de ter falado moi pouco con el) e para Gonzalo Cao, por seren dos estudantes máis brillantes que coñecín, unhas mentes incribles e as persoas máis novas en aumentar a miña paixón pola ciencia.









Máis alá do mundo académico, non podo esquecer aos meus grandes amigos Ángel, Patri e Chekia, que, malia non entenderen completamente o que eu facía, sempre estiveron aí para escoitar. A Lara, Darío (e a Tomás, o profesor-DJ que fixo que quixese estudar química, faltan docentes así...), Silva e Espi, grazas por compartiren as penurias do doutoramento e polos bos momentos de diversión. E como esquecer os \textit{Escape Rooms}, os partidos de baloncesto (xogar xuntos no Baltar, os campionatos $3\times3$), os xogos de mesa, as ceas no \textit{As de Picas}, ou a gran cantidade de cafés os sábados..., e en especial as longas noites xogando a \textit{Smite}, \textit{DBD} ou \textit{Baldur's Gate} con Samuel, Xoel, Berto, Nai, Martín e Marcos. Unha amizade que comezou na escola e que continúa a día de hoxe é un agasallo excepcional, e non podería ser máis feliz de vervos a todos ao meu carón. Sodes os mellores!



Canto á familia na que nacín, esta tese non tería sido posible sen o apoio inquebrantable dos meus pais, María del Carmen e José (Quérovos un mundo, non sabedes canto. A pesar de non estarmos de acordo en moitas cousas, por diferentes experiencias da vida, a vosa foi máis complicada... Grazas por darme unha vida tan feliz e por darme unha educación que me permitiu chegar a onde estou! A pesar de non terdes vós esas oportunidades, tanto é así, que chego a botar de menos traballar no CB con vós). Mesmo sen comprenderen totalmente os detalles da miña investigación, nunca deixaron de preguntar e de aprender sobre como funciona o sistema académico (a miña nai ata empezou a ver \textit{The Big Bang Theory}! O mellor que me deu a tese foi poder administrar o tempo para poder pasar máis con eles... a pesar de ter que traballar moitas fins de semana). Á miña avoa, Lola (cos seus paseos por Cabicastro), e a Vicente, que sempre agardaron que seguise os pasos de Jorge Mira e me fixese profesor. E aos meus irmáns, Noelia (unha rapaza cun corazón enorme, que non entendo como ten tan pouca confianza nela mesma, se é a mellor! Por moito máis tempo xuntos e seguir vendo como consegues os teus obxectivos, o mundo é teu!), Julio e Cris (por seren os mellores irmáns maiores do mundo: xogos, comidas, paseos, deporte, motos de xoguete, os partidos de Nadal... eses 16 anos de diferenza non foron un impedimento para que fósedes os mellores). E como non aos meus sobriños Laura (a futura estrela do balonmán), Sara (a Tiktoker), Mikel (o videoxogador) e Álex (o rapaz que non falla un acertixo de matemáticas!) e á familia Belga, en especial a Jordy e Rubén. Se tiven unha infancia marabillosa, foi grazas a vós. Tanto se me tocaba facer de irmán pequeno coma de maior, os partidos de fútbol en Montalvo, xogar en Dorrón con Adrián, os capítulos e series (en especial \textit{Miércoles} e \textit{My Hero Academia}) e as carreiras de \textit{MarioKart} na \textit{Wii} son lembranzas que atesouro con moito cariño. Aínda que a vida cambie, estou seguro de que están por vir experiencias aínda mellores.

Tamén desexo darlle as grazas a Beatriz Roca, Manuel e Chemita Carballedo por acollerme coma un máis da familia durante case unha década.

Finalmente, debo darlle as grazas á miña compañeira de vida, Coral. Unha rapaza que conquistou o meu corazón hai case 10 anos, cando me deixaba perder para poder tomar un xeado contigo, ou cando faltaba ás clases do venres para ir a Lugo unhas horas e pasar parte da tarde xuntos... Grazas por estares aí a cada paso do camiño, nas boas e nas malas. O teu apoio, confianza e paciencia foron a miña áncora. A vida é complicada, e aínda nos enfrontamos a momentos nos que estamos máis lonxe do que nos gustaría, pero estou seguro de que os superaremos. Sobrevivimos a un \textit{PIR}, a unha tese, a mestrados e a todas as miñas dúbidas sobre o futuro, mesmo cando ti vías as cousas máis claras ca min. Algúns dos mellores momentos da miña vida foron ao teu lado, especialmente a nosa viaxe a Nova York, e agardo vivir moitos máis, moitas máis viaxes, moitos máis balnearios, moitas máis ceas, máis partidas de bolos, e, en xeral, moita máis vida xuntos. Quérote.

E como dixo Mike Breen, cando Jalen Brunson abraza á súa familia ao ser campión da NBA (e MVP das Finais!) cos New York Knicks: \textit{Life is not perfect, but there are perfect moments. And there's one right there}. Os meus momentos perfectos son con vós. Grazas!
